\clearpage
\VerdierPhysicalPageAnchor{0158}%
\begingroup\small\setlength{\parskip}{0pt}\setlength{\emergencystretch}{1.5em}
\VerdierSemanticAnchor{ch3-3-2-5-kernel-triangle}
设 $(Z,X,Y',u,f'g,v)$ 为 $\mathcal D$ 中的特异三角。图式
\[
\begin{tikzcd}[column sep=5em,row sep=2.1em]
h(Z)\arrow[r,"h(fu)"]\arrow[d,"h(u)"']&h(Y)\arrow[d,"\operatorname{id}_{h(Y)}"]\\
h(X)\arrow[r,"h(f)"]\arrow[d,"h(g)"']&h(Y)\arrow[d,"h(g')"]\\
h(X')\arrow[r,"h(f')"']&h(Y')
\end{tikzcd}
\]
在水平方向取像后给出
\VerdierSemanticAnchor{ch3-3-2-5-1}
\[
\operatorname{Im}(fu)\xrightarrow{\beta}V\xrightarrow{\alpha}V'.
\tag{3.2.5.1}
\]
\VerdierSemanticAnchor{ch3-3-2-5-kernel-proof}
我们证明态射 $\operatorname{Im}(fu)\to V$ 是 $\alpha$ 的核。首先，
$\beta$ 是单态射：事实上，$\operatorname{Im}(fu)$ 与 $V$ 都是
$h(Y)$ 的子对象。其次，$\alpha\beta=0$：典范态射
$h(Z)\to\operatorname{Im}(fu)$ 是满态射，典范态射
$V'\to h(Y')$ 是单态射，而 $h(f'gu)$ 是 $h$ 作用于一个特异
三角中两个相继态射的复合所得的态射，故其为零 (1.2.2)。
还需证明在 $V$ 处正合。设 $W$ 为 $\mathcal D$ 的对象。任意态射
$\sigma\colon h(W)\to V$ 都来自 $h(W)$ 到 $h(X)$ 的态射，因为
$h(W)$ 在 $\mathcal D^\wedge$ 中是投射对象；也就是说，它来自一个态射
$m\colon W\to X$。它与 $\alpha$ 的复合为零，当且仅当
$f'gm=0$，因为 $V'$ 是 $h(Y')$ 的子对象。由 (1.2.1)，存在
$n\colon W\to Z$ 使 $un=m$。故 $\sigma$ 经 $\operatorname{Im}(fu)$ 分解。

\VerdierSemanticAnchor{ch3-3-2-5-cokernel-setup}
现证 $\alpha$ 在 $\mathcal D^\wedge$ 中的余核是 $\mathcal D^\wedge$ 的可容许对象。
设 $(X,Y',Z',g'f,u',v')$ 为 $\mathcal D$ 中的特异三角。图式
\[
\begin{tikzcd}[column sep=5em,row sep=2em]
h(X)\arrow[r,"h(f)"]\arrow[d,"h(g)"']&h(Y)\arrow[d,"h(g')"]\\
h(X')\arrow[r,"h(f')"]\arrow[d,"\operatorname{id}_{h(X')}"']&h(Y')\arrow[d,"h(u')"]\\
h(X')\arrow[r,"h(u'f')"']&h(Z')
\end{tikzcd}
\]
\endgroup
